\paragraph{Expression dispersion partly explains the cohort gap.}
\label{sec:cohortgap}
Holding tissue eligibility fixed at reference-only, the complete SMC and KUL3
mature-cell pools contain 662 and 844 cells. Across their 45 and 15 eligible
patient pairs, median source-pool squared coefficients of variation
($c_P^2=\sigma_P^2/\mu_P^2$) are 4.35 and 16.99; median zero-expression
fractions are 39.3\% and 62.1\%. The source means are 19.5 and 11.9 raw
counts, respectively (medians over pools). The KUL3 contrast therefore has
substantially more relative sampling noise despite its larger source-cell
pool.

For the implemented binomial thinning, the tumor-cell variance is
$s^2\sigma_P^2+s(1-s)\mu_P$. At $s=0.5$, the exact conditional variance of
the contrast relative to its squared signal is
\begin{equation}
\frac{\operatorname{Var}(\hat\theta\mid P)}{\theta_P(0.5)^2}
=\frac{c_P^2+1/\mu_P}{n_T}+\frac{4c_P^2}{n_N}.
\label{eq:relativevariance}
\end{equation}
With $n_T=50$ and $n_N=800$, plug-in signal-to-standard-deviation ratios
using the median pool moments are 3.02 in SMC and 1.53 in KUL3. This provides a distributional explanation
for weaker KUL3 discrimination. It does not reproduce the exact 70-versus-400
cutoff ratio: percentile-bootstrap behavior, pair heterogeneity and binning
also affect that selection. The biological causes of the source-distribution
difference remain unexplained.
